\documentclass[oneside,a4paper,14pt]{extarticle}
\usepackage[a4paper,letterpaper,top=20mm,bottom=20mm,left=20mm,right=10mm]{geometry}
\usepackage[russian]{babel}
\usepackage{indentfirst}
\usepackage{graphicx}
\usepackage{caption}
\usepackage{titlesec}
\usepackage{minted, fancyvrb}
\usepackage{hyperref}
\usepackage{enumitem}
\usepackage{amsmath}
\usepackage{fancyhdr}
\usepackage{colortbl}
\usepackage[dvipsnames]{xcolor}
\usepackage{float}

\setminted{style = rainbow_dash, fontsize = \small}

\titleformat{\section}{\normalsize\bfseries}{\thesection}{1em}{}
\titleformat{\subsection}{\normalsize\bfseries}{\thesubsection}{1em}{}
\titleformat{\subsubsection}{\normalsize\bfseries}{\thesubsubsection}{1em}{}

\renewcommand\baselinestretch{1.33}
\setlength{\parindent}{1.25cm}

\setlist[enumerate]{
  left=0.5\parindent,
  label=\arabic*.,
  itemsep=0pt,
  topsep=5pt,
  partopsep=0pt,
  parsep=0pt
}
\setlist[itemize]{
  left=0.5\parindent,
  itemsep=0pt,
  topsep=5pt,
  partopsep=0pt,
  parsep=0pt
}

\hypersetup{
  colorlinks=true,
  linkcolor=black,
  urlcolor=blue,
  pdfborder={0 0 0},
  pdftitle={Построение 2-разрядного сдвигового регистра на JK-триггерах},
  pdfauthor={Черкасов А.А. и Зинин В.А.},
}

\begin{document}

\newpage
\thispagestyle{empty}
\begin{center}
  МИНИСТЕРСТВО НАУКИ И ВЫСШЕГО ОБРАЗОВАНИЯ РОССИЙСКОЙ ФЕДЕРАЦИИ\\
  ФЕДЕРАЛЬНОЕ ГОСУДАРСТВЕННОЕ БЮДЖЕТНОЕ УЧРЕЖДЕНИЕ ВЫСШЕГО ОБРАЗОВАНИЯ\\
  <<ВЯТСКИЙ ГОСУДАРСТВЕННЫЙ УНИВЕРСИТЕТ>>\\
  Институт математики и информационных систем\\
  Факультет автоматики и вычислительной техники\\
  Кафедра электронных вычислительных машин
\end{center}
\vspace{10mm}

\hfill
\begin{tabular}{l}
  \footnotesize Дата сдачи на проверку:                                          \\
  \footnotesize <<\rule[-1mm]{5mm}{0.10mm}\/>>\rule[-1mm]{20mm}{0.10mm}\ 2025 г. \\
  \footnotesize Проверено:                                                       \\
  \footnotesize <<\rule[-1mm]{5mm}{0.10mm}\/>>\rule[-1mm]{20mm}{0.10mm}\ 2025 г. \\
\end{tabular}
\vfill

\begin{center}
  Отчёт по лабораторной работе №3\\
  по дисциплине\\
  <<Цифровая схемотехника>>\\
  \textbf{Построение 2-разрядного сдвигового регистра на JK-триггерах}\\
\end{center}
\vspace{25mm}
\noindent
\begin{tabular}{ll}
  Разработали студенты гр. ИВТб-2301-05-00 & \hspace{5mm}\rule[-1mm]{30mm}{0.10mm}\,/Черкасов А. А./ \\
                                           & \hspace{5mm}\rule[-1mm]{30mm}{0.10mm}\,/Зинин В. А./    \\
                                           & \hspace{12.5mm}\footnotesize(подпись)                   \\
  Преподаватель                           & \hspace{5mm}\rule[-1mm]{30mm}{0.10mm}\,/Мельцов В. Ю./  \\
                                           & \hspace{12.5mm}\footnotesize(подпись)                   \\
\end{tabular}

\noindent
\begin{tabular}{lp{58mm}r}
  Работа защищена &  & \hspace{5mm}<<\rule[-1mm]{5mm}{0.10mm}\/>>\rule[-1mm]{30mm}{0.10mm}\ 2026 г.
\end{tabular}
\vfill

\begin{center}
  Киров\\
  2026
\end{center}

\newpage\thispagestyle{plain}

\section{Цель лабораторной работы}
Получить навыки построения и анализа 2-разрядного сдвигового регистра на JK-триггерах. Изучить принцип последовательного сдвига данных. Выполнить синтез функциональной и принципиальной схем регистра на базе микросхем КП7, ТВ6 (серия ТТЛ).

\section{Теоретическое описание}
\subsection{JK-триггер}
JK-триггер является универсальным триггером, имеющим информационные входы $J$ и $K$, тактовый вход $C$ и асинхронные входы установки $\overline{S}$ и сброса $\overline{R}$. Вход $J$ отвечает за установку единицы, вход $K$ -- за сброс в ноль. При одновременной подаче единицы на оба входа ($J = K = 1$) триггер переключается в противоположное состояние (режим счёта).

В данной работе используются микросхемы КП7 (содержащая JK-триггеры) и ТВ6 (также JK-триггер с расширенными логическими входами).

\begin{center}
  Таблица 1 --- Таблица истинности JK-триггера \\
  \vspace{5mm}
  \begin{tabular}{|c|c|c|c|c|c|}
    \hline
    $\overline{S}$ & $\overline{R}$ & C & J & K & $Q_{n+1}$ \\
    \hline
    0 & 1 & X & X & X & 1 \\
    1 & 0 & X & X & X & 0 \\
    0 & 0 & X & X & X & 1* \\
    1 & 1 & $\downarrow$ & 0 & 0 & $Q_n$ \\
    1 & 1 & $\downarrow$ & 1 & 0 & 1 \\
    1 & 1 & $\downarrow$ & 0 & 1 & 0 \\
    1 & 1 & $\downarrow$ & 1 & 1 & $\overline{Q}_n$ \\
    \hline
  \end{tabular}
  \vspace{2mm}
  \par\footnotesize * --- неустойчивое состояние.
\end{center}

\subsection{Универсальный сдвиговый регистр}
Сдвиговый регистр -- последовательное устройство, в котором информация сдвигается от разряда к разряду под действием тактовых импульсов. В 2-разрядном универсальном сдвиговом регистре предусмотрено несколько режимов работы, задаваемых управляющими входами $S_0$ и $S_1$:
\begin{itemize}
  \item $S_1S_0 = 00$ --- хранение (состояние выходов не меняется);
  \item $S_1S_0 = 01$ --- сдвиг влево (последовательный вход $DL$);
  \item $S_1S_0 = 10$ --- сдвиг вправо (последовательный вход $DR$);
  \item $S_1S_0 = 11$ --- параллельная загрузка ($D_1$ --- старший разряд, $D_0$ --- младший).
\end{itemize}
Тактовые входы всех триггеров объединены и подключены к общему сигналу синхронизации $CLK$.

\begin{center}
  Таблица 2 --- Режимы работы сдвигового регистра \\
  \vspace{5mm}
  \begin{tabular}{|c|c|c|}
    \hline
    $S_1$ & $S_0$ & Режим \\
    \hline
    0 & 0 & Хранение \\
    0 & 1 & Сдвиг влево ($DL$) \\
    1 & 0 & Сдвиг вправо ($DR$) \\
    1 & 1 & Параллельная загрузка ($D_1, D_0$) \\
    \hline
  \end{tabular}
\end{center}

\section{Функциональная схема универсального сдвигового регистра}
На рисунке 1 представлена функциональная схема 2-разрядного универсального сдвигового регистра на JK-триггерах. Регистр имеет последовательные входы $DL$ (сдвиг влево) и $DR$ (сдвиг вправо), параллельные входы $D_1$ (старший разряд) и $D_0$ (младший разряд), управляющие входы $S_0$, $S_1$ и тактовый вход $CLK$. Выходы $Q_0$ и $Q_1$ отображают хранящееся значение. Комбинация сигналов на управляющих входах определяет режим работы регистра согласно таблице 2.

\begin{figure}[H]
  \centering
  \includegraphics[width=0.9\textwidth]{pics/functional_diagram.png}
  \caption*{Рисунок 1 --- Функциональная схема 2-разрядного сдвигового регистра на JK-триггерах}
\end{figure}

\section{Принципиальная схема устройства}
Принципиальная схема (рисунок 2) выполнена на микросхемах серии ТТЛ: \textbf{КП7} и \textbf{ТВ6}. КП7 содержит JK-триггеры, ТВ6 представляет собой JK-триггер с AND-входами на J и K. Из этих элементов собран 2-разрядный сдвиговый регистр. Схема включения микросхем соответствует заводским цоколёвкам. Питание подаётся от источника +5~В.

\begin{figure}[H]
  \centering
  \includegraphics[width=0.95\textwidth]{pics/principle_diagram.png}
  \caption*{Рисунок 2 --- Принципиальная схема 2-разрядного сдвигового регистра на JK-триггерах (КП7, ТВ6)}
\end{figure}

\section*{Вывод}
В ходе лабораторной работы были получены практические навыки построения 2-разрядного универсального сдвигового регистра на JK-триггерах. Регистр поддерживает четыре режима работы: хранение, сдвиг влево, сдвиг вправо и параллельную загрузку, управляемые сигналами $S_0$, $S_1$. Триггеры реализованы на базе микросхем КП7 и ТВ6 (серия ТТЛ). Разработаны функциональная и принципиальная схемы устройства. Полученные результаты могут быть использованы при проектировании более сложных сдвиговых регистров, преобразователей последовательного кода в параллельный и других узлов цифровой техники.

\end{document}
